\section{Задание}

В данной лабораторной работе моделируется работа информационно-вычислительного центра с произвольным числом операторов и накопителей. Система представляет собой двухуровневую структуру обслуживания заявок.

Центр состоит из N операторов и M накопителей (компьютеров для обработки данных). Каждый оператор может быть соединен с одним из накопителей. Заявки поступают в систему, обрабатываются операторами, после чего передаются на соответствующие накопители для дальнейшей обработки.

Генератор заявок характеризуется количеством генерируемых заявок, заданным пользователем. Интервал времени между заявками и его погрешность (среднее квадратичное отклонение) задаются вручную. Время поступления заявок определяется нормальным распределением.

Каждый из N операторов имеет свое время обработки заявки и погрешность. Время обработки определяется нормальным распределением. Каждый оператор связан с одним из M накопителей.

Каждый из M накопителей имеет фиксированное время обработки заявки. Накопитель может обслуживать заявки от нескольких операторов одновременно через механизм очереди.

Система работает по принципу событийного моделирования. Отказ заявки происходит в следующих случаях: если все N операторов заняты в момент поступления заявки (отказ на уровне оператора), либо если накопитель, к которому привязан оператор, обработавший заявку, занят обработкой другой заявки (отказ на уровне накопителя). Таким образом, даже если заявка была принята оператором, она может получить отказ при попытке передачи на накопитель.

Целью работы является проведение статистического анализа работы информационно-вычислительного центра с произвольной конфигурацией операторов и накопителей, определение вероятностей отказа на каждом уровне обслуживания и оценка общей эффективности системы.

\section{Теоретическая часть}

\subsection{Схемы модели}

На рисунке 1 представлена структурная схема модели с произвольной конфигурацией связей между операторами и компьютерами.

\includeimage
	{blockDiagram}{f}{H}{0.8\textwidth}
	{Структурная схема модели}

Разработанная модель поддерживает настраиваемую топологию связей между операторами и накопителями. Каждый оператор может направлять запросы на любой из компьютеров, что позволяет исследовать различные конфигурации информационного центра и их влияние на эффективность обработки запросов.

В процессе взаимодействия клиентов с информационным центром возможно два режима работы:

\begin{itemize}
    \item режим нормального обслуживания, когда клиент выбирает одного из свободных операторов, отдавая предпочтение тому, у кого максимальная производительность;
    \item режим отказа клиенту в обслуживании, когда все операторы заняты.
\end{itemize}

На рисунке 2 представлена схема модели в терминах систем массового обслуживания (СМО).

\includeimage
	{queuingSystems}{f}{H}{0.7\textwidth}
	{Схема модели в терминах СМО}

\subsection{Событийный принцип моделирования}

Событийный принцип использует дискретный характер изменения состояний моделируемой системы. Состояния блоков имитационной модели анализируются только в момент появления события. Момент наступления следующего события определяется минимальным значением из списка будущих событий, представляющего совокупность ближайших изменений состояния каждого блока системы.

Этот подход более эффективен по сравнению с пошаговым, так как устраняет необходимость проверки состояния системы на каждом временном шаге.

В разработанной модели определены следующие типы событий:

\begin{itemize}
    \item \textbf{CLIENT\_ARRIVAL} — прибытие клиента в информационный центр
    \item \textbf{OPERATOR\_FINISH} — окончание обработки запроса оператором
    \item \textbf{COMPUTER\_FINISH} — окончание обработки запроса компьютером
\end{itemize}

\subsection{Равномерное распределение}

Случайная величина $X$ имеет равномерное распределение на отрезке $[a, b]$, если её плотность распределения $f(x)$ равна:

\begin{equation}
p(x) = \begin{cases}
\frac{1}{b-a}, & \text{если } a \leq x \leq b \\
0, & \text{иначе}
\end{cases}
\end{equation}

При этом функция распределения $F(x)$ равна:

\begin{equation}
F(x) = \begin{cases}
0, & x < a \\
\frac{x-a}{b-a}, & a \leq x \leq b \\
1, & x > b
\end{cases}
\end{equation}

Обозначение: $X \sim R[a, b]$.

Интервал времени между появлением $i$-й и $(i-1)$-й заявки по равномерному закону вычисляется следующим образом:

\begin{equation}
T_i = a + (b - a) \cdot R
\end{equation}

где $R$ — псевдослучайное число из интервала $[0, 1]$.

\subsection{Переменные имитационной модели}

\textbf{Экзогенные переменные} (входные параметры модели):

\begin{itemize}
    \item параметры равномерного распределения для генератора клиентов ($a_{\text{gen}}$, $b_{\text{gen}}$);
    \item параметры равномерного распределения для каждого оператора ($a_i$, $b_i$);
    \item детерминированное время обработки для каждого компьютера ($t_j$);
    \item количество операторов ($N_{\text{op}}$);
    \item количество компьютеров ($N_{\text{comp}}$);
    \item топология связей между операторами и компьютерами;
    \item целевое количество клиентов для моделирования ($N_{\text{target}}$).
\end{itemize}

\textbf{Эндогенные переменные} (результаты моделирования):

\begin{itemize}
    \item конкретные времена обработки запросов (сгенерированные по распределениям);
    \item $n_0$ — число обслуженных клиентов;
    \item $n_1$ — число клиентов, получивших отказ;
    \item размеры очередей (текущий, максимальный, средний);
    \item состояния операторов и компьютеров в каждый момент времени.
\end{itemize}

Вероятность отказа в обслуживании клиента вычисляется как:

\begin{equation}
P = \frac{n_1}{n_0 + n_1}
\end{equation}

\section{Описание реализации}

\subsection{Используемые библиотеки}

Программа реализована на языке Python с использованием следующих библиотек:

\begin{itemize}
    \item \textbf{PyQt6} — создание графического интерфейса пользователя
    \item \textbf{random} — генерация псевдослучайных чисел
    \item \textbf{typing} — аннотации типов для улучшения читаемости кода
    \item \textbf{enum} — определение перечислений для типов событий
\end{itemize}

\subsection{Структура программы}

Программа состоит из следующих модулей:

\begin{itemize}
    \item \texttt{main.py} — точка входа в приложение
    \item \texttt{gui/main\_window.py} — главное окно с параметрами моделирования
    \item \texttt{gui/connections\_window.py} — окно настройки связей между операторами и компьютерами
    \item \texttt{gui/results\_window.py} — окно отображения результатов
    \item \texttt{models/distributions.py} — классы для генерации случайных величин
    \item \texttt{models/elements.py} — классы элементов системы (генератор, операторы, компьютеры, очереди)
    \item \texttt{models/event\_model.py} — реализация событийного подхода моделирования
    \item \texttt{constants.py} — константы интерфейса и параметры по умолчанию
\end{itemize}

\subsection{Метод решения}

Моделирование информационного центра выполняется с использованием событийного подхода:

\begin{enumerate}
    \item Инициализация системы:
    \begin{itemize}
        \item Создание генератора клиентов с заданным распределением
        \item Создание операторов с индивидуальными параметрами производительности
        \item Создание компьютеров с детерминированным временем обработки
        \item Создание очередей-накопителей для каждого компьютера
        \item Настройка связей между операторами и компьютерами
    \end{itemize}
    \item Добавление первого события прибытия клиента в список событий
    \item Основной цикл моделирования до обработки заданного числа клиентов:
    \begin{itemize}
        \item Выбор ближайшего события из отсортированного списка
        \item Обновление текущего времени моделирования
        \item Обработка события в зависимости от его типа:
        \begin{itemize}
            \item \textbf{Прибытие клиента}: поиск свободного оператора с максимальной производительностью или отказ в обслуживании
            \item \textbf{Окончание обработки оператором}: добавление запроса в очередь соответствующего компьютера (согласно топологии связей) и запуск обработки при возможности
            \item \textbf{Окончание обработки компьютером}: извлечение следующего запроса из очереди при наличии
        \end{itemize}
        \item Добавление новых событий в список с сохранением сортировки по времени
        \item Фиксация состояния очередей для расчёта статистики
    \end{itemize}
    \item Сбор финальной статистики по очередям (максимальный и средний размер)
\end{enumerate}

Производительность оператора вычисляется как величина, обратная среднему времени обработки:

\begin{equation}
\text{productivity}_i = \frac{1}{\frac{a_i + b_i}{2}} = \frac{2}{a_i + b_i}
\end{equation}

Клиент выбирает свободного оператора с наибольшей производительностью.

\subsection{Настройка топологии связей}

Программа позволяет настраивать связи между операторами и компьютерами через специальное графическое окно. Каждый оператор может быть связан с любым компьютером, что обеспечивает гибкость в исследовании различных конфигураций системы.

При обработке события окончания работы оператора запрос автоматически направляется в накопитель того компьютера, с которым данный оператор связан согласно заданной топологии.

\section{Практическая часть}

На рисунке 3 представлен результат работы разработанной программы для нахождения вероятности отказа. В качестве примера используется конфигурация с тремя операторами и двумя компьютерами:

\textbf{Параметры генератора:}
\begin{itemize}
    \item Равномерное распределение с параметрами $a = 8$, $b = 12$ (среднее 10 ± 2 минуты)
\end{itemize}

\textbf{Параметры операторов:}
\begin{itemize}
    \item Оператор 1: $a = 15$, $b = 25$ (среднее 20 ± 5 минут)
    \item Оператор 2: $a = 30$, $b = 50$ (среднее 40 ± 10 минут)
    \item Оператор 3: $a = 20$, $b = 60$ (среднее 40 ± 20 минут)
\end{itemize}

\textbf{Параметры компьютеров:}
\begin{itemize}
    \item Компьютер 1: детерминированное время обработки 15 минут
    \item Компьютер 2: детерминированное время обработки 30 минут
\end{itemize}

\textbf{Топология связей:}
\begin{itemize}
    \item Операторы 1 и 2 направляют запросы на Компьютер 1
    \item Оператор 3 направляет запросы на Компьютер 2
\end{itemize}

\textbf{Количество клиентов:} 300

\includeimage
	{result}{f}{H}{1.0\textwidth}
	{Результат работы программы}

Программа корректно вычисляет вероятность отказа в обслуживании клиентов. Событийный подход обеспечивает точное моделирование без погрешностей, связанных с дискретизацией времени.

Возможность настройки топологии позволяет проводить сравнительный анализ различных конфигураций системы и выбирать оптимальную схему распределения нагрузки между компьютерами.

\section{Вывод}

В ходе выполнения лабораторной работы была разработана программа с графическим интерфейсом для моделирования работы информационного центра с использованием событийного принципа и поддержкой настраиваемой топологии связей между операторами и накопителями.

Программа обеспечивает гибкую настройку параметров генератора клиентов, операторов и компьютеров, визуальную настройку топологии связей, автоматическую валидацию входных данных и наглядное представление результатов моделирования. Реализованный событийный подход обеспечивает точное моделирование без погрешностей дискретизации, характерных для пошагового метода.

Ключевая особенность разработанной программы — возможность динамического изменения количества операторов и компьютеров, а также настройки произвольных связей между ними. Это делает программу универсальным инструментом для исследования различных конфигураций систем массового обслуживания и анализа их эффективности.
